\documentclass[12pt]{article}
\usepackage[a4paper,margin=1in]{geometry}
\usepackage[T1]{fontenc}
\usepackage{lmodern,microtype}
\usepackage{amsmath,amssymb,amsthm,mathtools}
\usepackage{booktabs,array,enumitem,xcolor}
\usepackage[numbers,sort&compress]{natbib}
\usepackage[colorlinks=true,linkcolor=blue!55!black,citecolor=blue!55!black,urlcolor=blue!55!black]{hyperref}
\linespread{1.07}

\newtheorem{theorem}{Theorem}[section]
\newtheorem{proposition}[theorem]{Proposition}
\newtheorem{corollary}[theorem]{Corollary}
\theoremstyle{remark}
\newtheorem{remark}[theorem]{Remark}
\newcommand{\cH}{\mathcal H}
\newcommand{\norm}[1]{\left\lVert#1\right\rVert}

\title{Oblique Feshbach Determinant Factorization\\
\large Exact Packet--Complement Coordinates for a Biorthogonal Temporal Clock}
\author{Liang Wang}
\date{July 2026}

\begin{document}
\maketitle

\begin{abstract}
RH-188 identifies a surviving directional coupling product but leaves the
packet and complement resolvents undefined.  This paper supplies the exact
finite block decomposition for a biorthogonal temporal packet.

Let $V,W\in\mathbb C^{n\times r}$ satisfy $W^*V=I_r$.  Choose a frame $Z$
for $\ker W^*$ and set $S=[V,Z]$.  Then $S$ is invertible and its inverse has
the form $S^{-1}=[W^*;Y^*]$, with $Y^*V=0$ and $Y^*Z=I$.  In these coordinates
\[
 S^{-1}AS=
 \begin{pmatrix}K&B\\C&D\end{pmatrix},
 \quad
 K=W^*AV.
\]
Whenever $zI-D$ is invertible,
\[
 \det(zI-A)
 =\det(zI-D)
 \det\!\left(zI-K-B(zI-D)^{-1}C\right).
\]
This is a type-correct oblique Feshbach factorization.  The self-energy is the
exact object controlled by the directed Schur product.

A 240-case complex random audit verifies coordinate inversion, block
similarity, complement-gauge invariance, the one-factor norm bound, and the
determinant identity with zero failures.  The maximum relative determinant
error is $1.85\times10^{-11}$.  The theorem converts the
local bi-Krylov candidate into a precise next question: validate the
complement resolvent on root contours.  It does not itself supply those
bounds or prove a physical Riesz shell.
\end{abstract}

\section{Why an exact coordinate theorem is needed}

The RH-185 packet is not orthogonal.  Writing $Q=I-VW^*$ and informally
calling $QAQ$ the complement block hides two issues: $Q$ is oblique, and an
operator on $\operatorname{Ran}Q$ needs a chosen coordinate norm.  A
determinant identity should be derived from an actual similarity, not from a
symbolic orthogonal block matrix.

Finite-dimensional Feshbach and Schur complement identities are classical
\citep{GohbergGoldbergKaashoek1990}.  The contribution here is to state them
in the exact packet type produced by RH-184--185 and to make the complement
coordinates explicit.

\section{Biorthogonal packet coordinates}

Let $A\in\mathbb C^{n\times n}$ and let
$V,W\in\mathbb C^{n\times r}$ satisfy
\begin{equation}\label{eq:biorthogonal}
 W^*V=I_r.
\end{equation}
The associated oblique projector is $P=VW^*$.

Choose any full-rank
\begin{equation}\label{eq:Z}
 Z:\mathbb C^{n-r}\longrightarrow\ker W^*.
\end{equation}

\begin{proposition}[Invertible packet--complement coordinates]\label{prop:coordinates}
The matrix
\begin{equation}\label{eq:S}
 S=[V,Z]
\end{equation}
is invertible.  Its inverse has a unique block-row representation
\begin{equation}\label{eq:Sinv}
 S^{-1}=\begin{bmatrix}W^*\\Y^*\end{bmatrix}
\end{equation}
with
\begin{equation}\label{eq:dual-relations}
 Y^*V=0,
 \qquad Y^*Z=I_{n-r}.
\end{equation}
\end{proposition}

\begin{proof}
Suppose $Va+Zb=0$.  Applying $W^*$ and using
$W^*V=I$, $W^*Z=0$ gives $a=0$.  Then $Zb=0$, hence $b=0$ because $Z$ has
full rank.  Thus $S$ is invertible.  Multiplying $S^{-1}S=I$ shows that its
top row must be $W^*$ and its bottom row satisfies
\eqref{eq:dual-relations}.
\end{proof}

The choice of $Z$ is a complement-coordinate gauge.  Different choices are
related by an invertible block-triangular similarity and do not change the
full determinant.

For norm estimates there is a canonical convenient choice.

\begin{proposition}[Orthonormal complement coordinates]\label{prop:orthonormal-complement}
Let $Z$ have orthonormal columns spanning $\ker W^*$ and put
$Q=I-VW^*$.  Then the lower block row of $S^{-1}$ is
\begin{equation}\label{eq:Y-orthonormal}
 Y^*=Z^*Q.
\end{equation}
Consequently
\begin{equation}\label{eq:D-orthonormal}
 D=Z^*QAZ,
 \qquad
 \norm D\le\norm Q\norm A.
\end{equation}
\end{proposition}

\begin{proof}
Because $QZ=Z$ and $QV=0$, the row $Z^*Q$ satisfies
$(Z^*Q)V=0$ and $(Z^*Q)Z=I$.  Uniqueness in
Proposition~\ref{prop:coordinates} gives \eqref{eq:Y-orthonormal}.  The block
formula and its norm bound follow immediately.
\end{proof}

This removes one unnecessary oblique factor from the elementary complement
bound.  The ambient operator $QAQ$ obeys the valid but looser estimate
$\norm{QAQ}\le\norm Q^2\norm A$; the actual Feshbach coordinate block in an
orthonormal complement needs only one factor $\norm Q$.

\section{Exact block operator}

Define
\begin{equation}\label{eq:blocks}
 S^{-1}AS=
 \begin{pmatrix}
 K&B\\C&D
 \end{pmatrix},
\end{equation}
where
\begin{align}
 K&=W^*AV,
 &B&=W^*AZ,
 \label{eq:KB}\\
 C&=Y^*AV,
 &D&=Y^*AZ.
 \label{eq:CD}
\end{align}
The packet block $K$ is exactly the RH-185 compression.  The off-diagonal
blocks are coordinate forms of the RH-188 directed couplings.

\begin{proposition}[Residual interpretation]\label{prop:residual}
The ambient right residual satisfies
\begin{equation}\label{eq:right-residual}
 AV-VK=ZC,
\end{equation}
and the left residual satisfies
\begin{equation}\label{eq:left-residual}
 A^*W-WK^*=YB^*.
\end{equation}
\end{proposition}

\begin{proof}
From $AS=S\begin{psmallmatrix}K&B\\C&D\end{psmallmatrix}$, the first block
column gives $AV=VK+ZC$.  Taking the adjoint of the top block row of
$S^{-1}A=\begin{psmallmatrix}K&B\\C&D\end{psmallmatrix}S^{-1}$ gives the
second identity.
\end{proof}

Thus the numerical residuals and the Feshbach couplings are the same data in
different coordinate norms.

If $Z$ is replaced by $Z'=ZG$ for an invertible complement gauge $G$, then
$Y'^*=G^{-1}Y^*$ and
\begin{equation}\label{eq:complement-gauge}
 B'=BG,
 \qquad C'=G^{-1}C,
 \qquad D'=G^{-1}DG.
\end{equation}
It follows that
\begin{equation}\label{eq:complement-self-energy-invariance}
 B'(zI-D')^{-1}C'=B(zI-D)^{-1}C.
\end{equation}
Thus the self-energy is independent of the complement basis, even though
the three separate block norms depend on that basis.

\section{Feshbach determinant identity}

For a spectral parameter $z$ with $zI-D$ invertible, define the self-energy
\begin{equation}\label{eq:self-energy}
 \Sigma(z)=B(zI-D)^{-1}C
\end{equation}
and the reduced Feshbach matrix
\begin{equation}\label{eq:feshbach-matrix}
 F(z)=zI-K-\Sigma(z).
\end{equation}

\begin{theorem}[Oblique determinant factorization]\label{thm:determinant}
Whenever $zI-D$ is invertible,
\begin{equation}\label{eq:determinant}
 \boxed{
 \det(zI-A)=\det(zI-D)\det F(z).
 }
\end{equation}
\end{theorem}

\begin{proof}
Similarity by $S$ gives
\[
 \det(zI-A)
 =\det\begin{pmatrix}zI-K&-B\\-C&zI-D\end{pmatrix}.
\]
Multiply on the left by the determinant-one block matrix
\[
 \begin{pmatrix}I&B(zI-D)^{-1}\\0&I\end{pmatrix}.
\]
The result is block lower triangular with diagonal blocks
$F(z)$ and $zI-D$.  Taking determinants proves
\eqref{eq:determinant}.
\end{proof}

\begin{corollary}[Schur sufficient condition]\label{cor:schur}
If $zI-K$ and $zI-D$ are invertible and
\begin{equation}\label{eq:schur-condition}
 \norm{(zI-K)^{-1}}
 \norm B
 \norm{(zI-D)^{-1}}
 \norm C<1,
\end{equation}
then $F(z)$ is invertible and $z\notin\sigma(A)$.
\end{corollary}

\begin{proof}
Factor
$F(z)=(zI-K)[I-(zI-K)^{-1}\Sigma(z)]$ and apply the Neumann lemma.
\end{proof}

This is the exact location of the four factors recorded architecturally in
RH-163 and numerically separated in RH-188.

\section{Riesz rank consequence}

The determinant identity also gives the precise counting statement.  The
complement count cannot be silently omitted.

\begin{theorem}[Riesz count under a uniform Schur margin]\label{thm:riesz-count}
Let $\Gamma$ be a positively oriented simple contour on which both
$zI-K$ and $zI-D$ are invertible.  Suppose
\begin{equation}\label{eq:uniform-schur}
 \sup_{z\in\Gamma}
 \norm{(zI-K)^{-1}}\norm B
 \norm{(zI-D)^{-1}}\norm C<1.
\end{equation}
Then the algebraic eigenvalue count satisfies
\begin{equation}\label{eq:count-sum}
 N_A(\Gamma)=N_K(\Gamma)+N_D(\Gamma).
\end{equation}
Here $N_T(\Gamma)$ denotes the algebraic multiplicity of the spectrum of
$T$ inside $\Gamma$.
In particular, $N_A(\Gamma)=N_K(\Gamma)$ if the complement block has no
eigenvalues inside $\Gamma$.
\end{theorem}

\begin{proof}
For $0\le t\le1$, define
\[
 A_t=S\begin{pmatrix}K&tB\\tC&D\end{pmatrix}S^{-1}.
\]
Its reduced matrix is
$F_t(z)=zI-K-t^2B(zI-D)^{-1}C$.  The bound
\eqref{eq:uniform-schur} makes $F_t(z)$ invertible on $\Gamma$ for every
$t$.  The full determinant therefore has no contour zero during the
homotopy, so its winding number is constant.  At $t=0$ the operator is block
diagonal and its count is $N_K+N_D$.
\end{proof}

The theorem is conditional on continuous packet and complement inverse
bounds.  To isolate one packet root, one must either validate
$N_D(\Gamma)=0$ or compute and retain the complement contribution.  Small
residuals alone imply neither fact.

\section{Finite identity audit}

The audit uses dimensions $5,7,9,11$, packet ranks $1,2,3$, and twenty
complex random trials per pair, for 240 cases.  The left frame is corrected
to satisfy $W^*V=I$ exactly at floating precision.  A numerical nullspace of
$W^*$ supplies $Z$, and $Y$ is read from $S^{-1}$.

The maximum errors are:
\begin{center}
\begin{tabular}{lr}
\toprule
identity & maximum error\\
\midrule
$S^{-1}S=I$ & $3.22\times10^{-14}$\\
block similarity & $1.01\times10^{-11}$\\
relative determinant factorization & $1.85\times10^{-11}$\\
orthonormal complement dual row & $2.38\times10^{-12}$\\
one-factor norm-bound violation & $0$\\
complement-gauge self-energy & $2.64\times10^{-12}$\\
\bottomrule
\end{tabular}
\end{center}
All 240 cases pass the declared $10^{-8}$ tolerance.

The identity audit includes nominal recomputation of the self-energy in a
second complement gauge.  A later validated implementation must additionally
enclose that equality.  This catches a numerically important failure mode:
separate block norms can change greatly under a poor complement basis even
though the exact reduced determinant is fixed.

For outward work, the orthonormal choice in
Proposition~\ref{prop:orthonormal-complement} is the baseline because it
gives the one-factor bound $\norm D\le\norm Q\norm A$.  Another gauge is
useful only if the transformation and its improved product bound are both
validated.

\section{Physical next step}

For each surviving RH-185 window, the theorem defines a finite exact program:
\begin{enumerate}[leftmargin=2em]
\item build outward enclosures of $V,W,K,B,C,D$;
\item choose root contours from the predeclared length-four geometry;
\item validate $(zI-D)^{-1}$ on a finite contour mesh;
\item combine mesh, operator-ball, and coupling errors using the RH-167--168
 machinery summarized in RH-171 \citep{WangRH171};
\item apply \eqref{eq:schur-condition}, compute $N_D$, and use
 Theorem~\ref{thm:riesz-count} to count the full Riesz rank.
\end{enumerate}
The only genuinely new hard object is the physical complement resolvent.
The packet and coupling types are now exact.

\section{Boundary}

This paper proves a finite determinant identity for every biorthogonal packet
pair.  It does not validate any physical complement inverse, establish a
continuous Schur margin, prove a physical Riesz shell, transport shells
across scales, close R or Gate A, or approach a self-adjoint Hilbert--Polya
operator or the Riemann Hypothesis.

\bibliographystyle{plainnat}
\bibliography{references}
\end{document}
